\chapter{Introduction}
\label{ch:intro}

\section{What is MFTL2?}

\mftl{} is a C99 library for triangulated surfaces that are distributed over the ranks of an MPI
program. It stores the surfaces, moves them between ranks as they move through space, computes
their geometry (normals, curvature, areas, volumes, centres, moments of inertia), answers
point and intersection queries, and keeps their mesh quality under control by remeshing. It is
the surface engine of the front-tracking and immersed-boundary methods of the \multi{} flow
solver, where a surface is advected by the flow, carries forces and scalars, and is
remeshed every few time steps.

The surfaces it handles are
\begin{itemize}
  \item \textbf{closed} surfaces (bubbles, drops, particles) in a bounded or periodic box;
  \item \textbf{periodic} surfaces: a closed body that crosses a periodic boundary (a
    \emph{seam}), and surfaces that \emph{wrap} a periodic direction --- a ribbon that spans the
    box along $x$, a film that spans $x$ and $y$, a surface that wraps all three directions;
  \item \textbf{open} surfaces with a physical boundary (a membrane pinned on a ring, a cap on a
    wall orifice, a patch), declared as open by the application.
\end{itemize}
Every combination is supported: an open surface may cross periodic seams, a wrapping film may
be cut by many rank boundaries.

This manual describes \mftl{} as it exists after Stage~16 of its development (generalized
periodic and open surfaces) with the source layout of Stage~17. It explains the concepts
first --- the data structures, the parallel model, the invariants the library maintains ---
and then the functions that act on them. Every code listing in the manual is a complete
program from the \texttt{examples/} directory; the test suite compiles and runs them on 1, 2,
4 and 8 ranks, so they work as printed.

\section{Differences from MFTL1}
\label{sec:mftl1}

\mftl{} is a rewrite of the original MultiFlow Triangle Library (\mftlone). It keeps the
numerical conventions of \mftlone{} wherever \multi{} depends on them --- the mean-curvature
convention $2/R$, the remeshing criteria (with an explicit \mftlone{}-parity mode), the volume
recovery, the overlap quantification --- but replaces the data structures and the parallel
machinery. Table~\ref{tab:mftl1} summarises the differences. \mftlone{} identifiers appear in
this manual only in such historical comparisons.

\begin{table}[ht]
\centering\small
\begin{tabularx}{\textwidth}{@{}lXX@{}}
\toprule
 & \mftlone & \mftl \\
\midrule
Mesh storage & three hash tables (uthash) of vertices, edges and triangles & flattened arrays of vertices, triangles and half-edges per object; \texttt{khash} maps from global id to array index \\
Identity & sequential \texttt{id} & \texttt{global\_id} (unique per object across ranks) plus a rank-local array index \\
Ownership & \texttt{processor} field & \texttt{owner\_proc} and \texttt{is\_ghost}; exactly one owner per entity \\
Topology & edge records & half-edges with a twin state (LOCAL / REMOTE / BOUNDARY) \\
Decomposition & x-only 1-D ownership ranges & MPI Cartesian (or rectilinear) decomposition with a 26-neighbour graph communicator; the x-only scheme remains for load-balanced remeshing \\
Spatial search & oriented bounding-box tree & axis-aligned boxes and a linear bounding-volume hierarchy (LBVH) \\
Periodicity & not supported & periodic axes, image counters, winding classes, wrapping surfaces \\
Open surfaces & not supported & declared open bodies with tagged physical boundaries and boundary loops \\
Closed-phase quantities of general surfaces & --- & phase-boundary sets with an explicit DEFINED / UNDEFINED / UNSUPPORTED status \\
LAPACK & required & optional (paraboloid fits; a barycentric fit replaces it when absent) \\
\bottomrule
\end{tabularx}
\caption{\mftlone{} and \mftl{} compared.}
\label{tab:mftl1}
\end{table}

\section{Conventions used throughout}
\label{sec:conventions}

\begin{description}
\item[Return codes.] Functions that can fail return \texttt{MFTL\_SUCCESS} (0) or
  \texttt{MFTL\_ERROR} ($-1$). A collective function returns the same code on every rank.
\item[Collective vs.\ rank-local.] A \emph{collective} function must be called by every rank of
  the ensemble's communicator, in the same order and with the same arguments where the
  documentation says so (\eg the same query point). A \emph{rank-local} function works on the
  entities stored on the calling rank only. The documentation of every function in the headers
  states which it is.
\item[Winding.] Closed bodies are wound \emph{outward}: the triangle normal
  $(\vect{x}_1-\vect{x}_0)\times(\vect{x}_2-\vect{x}_0)$ points out of the body, and the
  enclosed volume is positive. Generators and readers canonicalise closed bodies to this
  winding.
\item[Curvature.] The stored mean curvature is $2H$, the magnitude of the discrete
  Laplace--Beltrami vector: $2/R$ on a sphere of radius $R$. The Gaussian curvature follows the
  \mftlone{} convention $2\times$(angle deficit)/(mixed area): $2/R^2$ on a sphere.
\item[Tolerances.] \texttt{MFTL\_VERYSMALL} ($10^{-12}$) and \texttt{MFTL\_VERYVERYSMALL}
  ($10^{-15}$) are fixed absolute guards against division by zero and degenerate geometry;
  \texttt{MFTL\_EPSILON} is $10^{-9}$. Decisions that depend on the scale of the mesh use
  mesh-relative tolerances inside the modules that take them. The overlap module uses the
  \mftlone{} slack \texttt{MFTL\_REL\_TOLERANCE} ($5\times10^{-7}$).
\item[Rank independence.] Integrals, loops, phase-set statuses and remeshing decisions that are
  documented as collective produce the same result on every rank. Integrals are reductions and
  may differ in the last bits between runs on different rank counts; the tests compare them with
  tolerances where that is the case.
\end{description}

\section{Glossary}
\label{sec:glossary}

\begin{description}[style=nextline,leftmargin=2.5em]
\item[object (body)] One surface, possibly with several connected components:
  \texttt{MFTL\_Object}. Identified by its object id (\texttt{global\_id} of the object).
\item[ensemble] The library handle \texttt{MFTL\_Ensemble}: one domain and the objects living
  in it. Every rank holds the same list of objects in the same order.
\item[domain] The global box, its periodicity and its decomposition over the ranks
  (\texttt{MFTL\_Domain}).
\item[slab] The box of one rank in the decomposition (the name comes from the x-only layout,
  where the boxes are slabs).
\item[vertex, triangle, half-edge] The entities of a mesh. A half-edge is a directed edge of
  exactly one triangle.
\item[twin] The opposite half-edge of the neighbouring triangle across an edge.
\item[global id / local index] A global id identifies an entity within its object on every
  rank; a local index is its position in this rank's array and changes with every
  redistribution.
\item[owner, ghost] Every entity has exactly one owner rank, whose copy is authoritative
  (\texttt{owner\_proc}). Copies on other ranks are read-only ghosts (\texttt{is\_ghost}).
\item[rim (MPI rim)] An edge whose twin is stored on another rank only: the edge of this
  rank's part of the mesh. A rim is \emph{not} a boundary.
\item[physical boundary] An edge with a single triangle on a declared-open body. It carries a
  non-negative integer \emph{tag}.
\item[boundary loop] A closed chain of boundary edges.
\item[seam] The crossing of a periodic face of the box by the surface. A seam is geometric, not a
  mesh attribute: the edges across it are ordinary interior edges.
\item[image counter] The integer vector \texttt{img} of a vertex that records which periodic
  image of the box the vertex belongs to.
\item[continuous coordinate] $\vect{X}=\texttt{pos}+\texttt{img}\odot\vect{L}$: the position of a
  vertex on the unbroken surface.
\item[minimum image] The periodic image of a difference vector that is shortest; exact for every
  mesh edge because edges are shorter than half a period.
\item[cut edge] An edge whose lattice shift is non-zero along some axis (\S\ref{sec:winding}).
\item[wrapping axis, compact axis] A periodic axis along which the body winds (the body spans the
  period), resp.\ does not.
\item[phase-boundary set] A list of bodies with orientation signs (and wall closures) that
  together bound a phase region (Chapter~\ref{ch:phasesets}).
\item[quiescent point] A moment between library operations at which all invariants of
  Chapter~\ref{ch:invariants} hold, \eg on return from a redistribution.
\end{description}

\section{Structure of this manual}

Chapter~\ref{ch:building} explains how to build, link and test the library.
Chapters~\ref{ch:datamodel}--\ref{ch:topology} introduce the data model, the parallel model and
the topological invariant. Chapters~\ref{ch:periodicity} and~\ref{ch:phasesets} cover periodicity
and the closed-phase quantities of general surfaces. Chapters~\ref{ch:geometry}--\ref{ch:data}
describe the geometric functions, generators and I/O, queries, remeshing and per-triangle data.
Chapter~\ref{ch:multiflow} explains how \multi{} uses the library, Chapter~\ref{ch:invariants}
collects the invariants, Chapter~\ref{ch:examples} the worked examples, and
Chapter~\ref{ch:limitations} the known limitations. Appendix~\ref{app:api} lists the public
API by header.
